\documentclass[11pt]{article}

\usepackage[margin=1in]{geometry}
\usepackage{amsmath,amssymb}
\usepackage[hidelinks]{hyperref}

\newcommand{\OmegaT}{\Omega_\theta}

\title{Literature-Implied Partial Subcase for arXiv:1701.07084}
\author{FA/Banach run fa\_banach\_001}
\date{June 14, 2026}

\begin{document}
\maketitle

\section*{Status}

This is a \texttt{literature\_implied\_answer (partial subcase)} packet. It
does not claim an original proof or a full answer to Correa's final question
for arbitrary Banach-space interpolation pairs. It records that a later
splitting-stability theorem of Castillo--Correa--Ferenczi--Gonzalez
\cite{CCFG2017} gives an affirmative answer for interpolation pairs of
superreflexive Kothe function spaces.

\section*{Source Question}

In Section 5, Final Remarks, Correa \cite{Correa2017} asks:
\[
  \text{If } dX_\theta \text{ is nontrivial for every }
  \theta\in(0,1),\ \theta\ne \frac12,\text{ is } dX_{1/2}
  \text{ nontrivial?}
\]
In the source paper, \(dX_\theta\) denotes the derived twisted sum associated
with the complex interpolation derivation \(\Omega_\theta\).

\section*{Supporting Theorem}

The decisive supporting result is Theorem 4.3 of
\cite{CCFG2017}, on PDF page 16:
if \((X_0,X_1)\) is an interpolation pair of superreflexive Kothe function
spaces and \(0<\theta<1\), then \(\Omega_\theta\) is trivial if and only if
there is a weight function \(w\) such that
\[
  X_1 = X_0(w)
\]
up to equivalence of norms.

\section*{Identification}

The endpoint condition \(X_1=X_0(w)\) is independent of \(\theta\). Therefore
Theorem 4.3 has the following immediate contrapositive consequence:

\medskip
\noindent
If \((X_0,X_1)\) is a superreflexive Kothe interpolation pair and
\(dX_\theta\) is nontrivial for at least one \(0<\theta<1\), then the endpoint
pair is not merely a weighted version of one Kothe space; hence no
\(\Omega_t\), \(0<t<1\), is trivial.

\medskip
\noindent
In particular, if \(dX_\theta\) is nontrivial for every
\(\theta\ne 1/2\), then \(dX_{1/2}\) is nontrivial.

\section*{Scope and Limitations}

This answers the source question only for interpolation pairs of
superreflexive Kothe function spaces. That includes the standard
\(\ell_p\)- and \(L_p\)-type Kothe settings motivating the question, but it
does not settle the arbitrary Banach-space pair case. The supporting paper
explicitly notes that the existence of local or global stability for the
differential process associated to complex interpolation of a general couple
of Banach spaces remains open.

The relation to Correa's final question is an agent-identified implication.
The supporting paper cites Correa's article and develops the relevant
splitting-stability theory, but it does not single out the final midpoint
question as the named problem being answered.

\section*{Search Bounds}

Before packaging, the run's cheap indexes were searched for
\texttt{1701.07084}, the paper title, \texttt{dX\_\{1/2\}},
\texttt{nontrivial for every theta}, and core twisted-sum/interpolation
keywords. Local later-source searches for citations to Correa's paper found
arXiv:1712.09647 as the decisive supporting theorem. No prior packet or
attempt for this exact Kothe subcase was found.

\section*{Human Review Recommendation}

Review as a provenance/status packet. The main points to check are the
translation between \(dX_\theta\) and the derivation-generated exact sequence
for \(\Omega_\theta\), and that the packet is restricted to superreflexive
Kothe function-space interpolation pairs.

\begin{thebibliography}{9}

\bibitem{Correa2017}
W. H. G. Correa,
\emph{Type, cotype and twisted sums induced by complex interpolation},
arXiv:1701.07084.

\bibitem{CCFG2017}
J. M. F. Castillo, W. H. G. Correa, V. Ferenczi, and M. Gonzalez,
\emph{Stability of the differential process generated by complex interpolation},
arXiv:1712.09647.

\end{thebibliography}

\end{document}
